% Pursuing Stacks (cotes 134-2 to 134-8): archive cote 134-5, site page 68 (PDF page 69), typescript page 255.
% typescript. Draft compared: ps4.tex:3425-3498.
% AI-transcribed by Claude Opus 5.5 (claude-opus-5-5, adaptive thinking, effort high), Message Batches API, batch ps2 (msgbatch_01NVwr1Yuo8qJf5rykeGBXPK), 2026-11-12; not yet checked by eye.
\scanpage{134-5/68}
\mspage{255}

\uline{Remark} 3. Assume \xout{that} in theorem 1 that A is totally aspheric. Then
\xout{the condition (ii} condition (ii) just means that i is $M_{as}$-aspheric, and
condition (iii) that $i^{*}(x)$ is aspheric for any x in C (as \xout{\illegible{1}} we got
$A^{\wedge}_{loc.as} = A^{\wedge}_{as} = A^{\wedge}_{tot.as}$ ). If moreover C satisfies the condition of
cor.\over{2}{3} above, then these conditions imply that \xout{A} A is a strict test cate\xout{g}%
gory, \xout{hence} and if $(M,M_{as})$ is modelizing, that i is a test functor
as stated in corollary 2. These observations sum up the substance of
the restatement of the main result of section 51, \inserted{given} in theorem 2 of p.178\xout{\illegible{1}}\hand{,}
\margin{\unsure{any one conv.?} p. 347 for \unsure{strong version}}in the present general context.

80. In the preceding section, we associated to any contractibility
structure $M_{c}$ on a category M, an asphericity structure $M_{as}$ ``generated''
by the former in a natural sense. It is this possibility of associating
(in a topologically meaningful way) an asphericity structure to a given
contractibility structure, which singles out the latter structure type,
among the three essential distinct``homotopy flavored'' kind of structures
developped at length in sections 51 \deleted{\illegible{1}}\inserted{and} 52 , in preference to the two
weaker notions of \inserted{a} homotopy interval structure\hand{,} and of a homotopism struc-%
ture (or, equivalently, of a \xout{\illegible{2}} \hand{``}homotopy relation\hand{''}). The
association
\[ (*) \qquad M_{c} \longmapsto M_{as} \quad \hand{(}\hand{$\longleftrightarrow$} \inserted{corresponding} notion of weak equivalence $W_{\hand{a}}$ in M\hand{)} \]
finally carried through \xout{here} in the last section, \over{has}{had} been forshadowed
\margin{110, \unsure{135}}earlier (cf. \inserted{p. 110, \deleted{135} and} p.142), but was pushed off for quite a while, in order
to``give precedence'' to the other approach in view \inserted{by then} towards more general
test functors\inserted{than before}\hand{,} leading up finally to the ``awkward main result'' of section
65. In \xout{\illegible{1}} (*) , as the asphericity structure $M_{as}$ is totally aspheric
and a fortiori aspheric, $M_{as}$ and the corresponding notion of weak equi-%
valence $W_{a}= W_{M_{c}}$ (which of course should not be confused with the notion
of homotopism\xout{s} associated to $M_{c}$ , giving a considerably smaller set of
arrows $W_{c}$ ) determine each other mutually. Till the writing up of secti
\hand{-}on 51, it was the aspect ``weak equivalence'' $W_{a}$ which was in the fore,
whereas the \xout{technically} conceptually more relevant aspect of ``aspheric
objects'' did not \xout{\illegible{6}}\inserted[typed]{appear in full light} before I was through with
grinding out the ``awkward approach'' (cf. p.188).

\deleted{Here} \hand{]}it is time now to remember the opposite association
\[ (\,**\,) \qquad W_{a} \longmapsto \text{homotopy structure } h_{W_{a}} \qquad , \]

\margin{68}
